\section*{Supplementary Material}

\setcounter{table}{0}
\renewcommand{\thetable}{S\arabic{table}}

\begin{table}[h!]
  \centering
  \caption{\textbf{The Biological Equivalence Principle.} Correspondence between the geometric variables of General Relativity (GR) and the morphoelastic variables of the Information-Energy-Countercurvature (IEC) framework. This analogy motivates the mathematical structure of the cost functional in Eq.~\ref{eq:cost_functional}.}
  \label{tab:gr_analogy}
  \small
  \begin{tabular}{p{3cm} p{5cm} p{5cm}}
    \toprule
    \textbf{Concept} & \textbf{General Relativity (GR)} & \textbf{IEC Framework} \\
    \midrule
    \textit{Geometry} & Metric tensor $g_{\mu\nu}$ & Reference configuration $\hat{\mathbf{u}}(s)$ \\
    \textit{Source} & Stress-Energy tensor $T_{\mu\nu}$ & Mechanical load vector $\mathbf{n}(s)$ \\
    \textit{Equation} & $R_{\mu\nu} - \frac{1}{2}Rg_{\mu\nu} = \kappa T_{\mu\nu}$ & $\mathbf{n}' + \mathbf{f} = \mathbf{0}, \mathbf{m}' + \mathbf{u}' \times \mathbf{n} = \mathbf{0}$ \\
    \textit{Path} & Geodesic ($\delta \int ds = 0$) & Homeostatic shape ($\delta \mathcal{E} = 0$) \\
    \textit{Information} & Information Term $I_{\mu\nu}$ & Information Field $I(s)$ \\
    \textit{Function} & Modulates space curvature & Modulates stress-free strain $\hat{\kappa}$ \\
    \bottomrule
  \end{tabular}
\end{table}
